%%==============================================================
%% TRABAJO UNIVERSITARIO — FORMATO APA 7 (sin clase apa7)
%% Todos los parámetros editables están en la sección
%% "CONFIGURACIÓN DEL DOCUMENTO" más abajo.
%%==============================================================
\documentclass[12pt]{article}

%% ── Codificación e idioma ─────────────────────────────────
\usepackage[utf8]{inputenc}
\usepackage[T1]{fontenc}
\usepackage[spanish]{babel}

%% ── Márgenes APA 7 (1 in por todos lados) ─────────────────
\usepackage[top=1in, bottom=1in, left=1in, right=1in]{geometry}

%% ── Interlineado doble ────────────────────────────────────
\usepackage{setspace}
\doublespacing

%% ── Fuente Times New Roman ────────────────────────────────
\usepackage{mathptmx}

%% ── Encabezado de página ──────────────────────────────────
\usepackage{fancyhdr}
\pagestyle{fancy}
\fancyhf{}
\renewcommand{\headrulewidth}{0pt}
\fancyhead[R]{\thepage}          % número de página arriba a la derecha

%% ── Justificación y sangría ───────────────────────────────
\usepackage{ragged2e}
\usepackage{indentfirst}

%% ── Secciones numeradas APA 7 ─────────────────────────────
%%   Nivel 1: centrado, negrita
%%   Nivel 2: izquierda, negrita
%%   Nivel 3: izquierda, negrita + cursiva
\usepackage{titlesec}
\setcounter{secnumdepth}{3}

\titleformat{\section}[block]
  {\normalfont\normalsize\bfseries}{\thesection.}{0.5em}{}
\titlespacing*{\section}{0pt}{12pt}{6pt}

\titleformat{\subsection}[block]
  {\normalfont\normalsize\bfseries}{\thesubsection.}{0.5em}{}
\titlespacing*{\subsection}{0pt}{12pt}{6pt}

\titleformat{\subsubsection}[block]
  {\normalfont\normalsize\bfseries\itshape}{\thesubsubsection.}{0.5em}{}
\titlespacing*{\subsubsection}{0pt}{12pt}{6pt}

%% ── Citas y bibliografía ──────────────────────────────────
\usepackage{csquotes}
\usepackage[style=apa, backend=biber]{biblatex}
\addbibresource{bibliography.bib}

%% ── Tablas y figuras ──────────────────────────────────────
\usepackage{booktabs}
\usepackage{threeparttable}
\usepackage{caption}
\captionsetup[table]{
  position=above,
  labelfont=bf,
  labelsep=newline,
  textfont=it,
  justification=raggedright,
  singlelinecheck=false
}
\captionsetup[figure]{
  position=above,
  labelfont=bf,
  labelsep=newline,
  textfont=it,
  justification=raggedright,
  singlelinecheck=false
}

%% ── Hipervínculos (sin color) ─────────────────────────────
\usepackage[hidelinks]{hyperref}

%% ── Texto de relleno (quitar en versión final) ────────────
\usepackage{lipsum}


%%==============================================================
%%  CONFIGURACIÓN DEL DOCUMENTO — edita solo esta sección
%%==============================================================

\newcommand{\docTitulo}{Brecha Digital en el acceso a las TIC en unidades
  educativas en el área rural}
\newcommand{\docTituloCorto}{Brecha digital en área rural}   % encabezado

\newcommand{\docAutores}{%
  Teresa Añez\textsuperscript{1}, 
  Mery Canza\textsuperscript{2}, 
  Olivia Bernal\textsuperscript{3}, 
  Neiberth Espinoza\textsuperscript{4},\\
  Cinthia Sanchez\textsuperscript{5}, 
  Fernando Miranda\textsuperscript{6}, 
  Luis Cartagena\textsuperscript{7}}

\newcommand{\docInstitucion}{Universidad Gabriel René Moreno}
\newcommand{\docAsignatura}{Fundamentos de la Pedagogía Crítica y Contemporánea}
\newcommand{\docDocente}{MSc. Lisbeth Katherine Morales Camacho}
\newcommand{\docFecha}{Sábado 4 de Julio de 2026}

\newcommand{\docResumen}{%
  El siguiente documento investiga las causas que provocan la brecha digital
  y las consecuencias que generan en los estudiantes de comunidades rurales de
  Bolivia. De igual manera se presentan acciones que deberían tomarse para
  mitigar este problema, ya sea mejorando la infraestructura actual como la
  capacitación del personal docente.}

\newcommand{\docPalabrasClave}{brecha digital, TIC, educación rural, Bolivia}

\newcommand{\docNotaAutor}{%
  Trabajo grupal correspondiente a la asignatura de \docAsignatura, Grupo~6.}

%%==============================================================


%%--------------------------------------------------------------
%% Encabezado con título corto (activo desde pág. 2)
%%--------------------------------------------------------------
\fancyhead[L]{\small\MakeUppercase{\docTituloCorto}}

%% Página 1: sin encabezado lateral, solo número
\fancypagestyle{portada}{%
  \fancyhf{}
  \renewcommand{\headrulewidth}{0pt}
  \fancyhead[R]{\thepage}
}

%%==============================================================
\begin{document}

%%--------------------------------------------------------------
%% PÁGINA 1 — PORTADA
%%--------------------------------------------------------------
\thispagestyle{portada}
\begin{center}
  \vspace*{2in}
  {\bfseries\docTitulo}\par
  \vspace{24pt}
  \docAutores\par
  \vspace{6pt}
  \docInstitucion\par
  \vspace{6pt}
  \docAsignatura\par
  \docDocente\par
  \docFecha\par
\end{center}

%% Nota de autor al pie de la portada
\vfill
\noindent\rule{\linewidth}{0.4pt}\par
\vspace{4pt}
\noindent{\bfseries Nota de Autor}\par
\vspace{4pt}
\noindent\docNotaAutor\par

\newpage

%%--------------------------------------------------------------
%% PÁGINA 2 — RESUMEN
%%--------------------------------------------------------------
\begin{center}
  {\bfseries Resumen}
\end{center}
\noindent\docResumen\par
\vspace{12pt}
\noindent\textit{Palabras clave:} \docPalabrasClave

\newpage

%%--------------------------------------------------------------
%% CUERPO DEL DOCUMENTO
%%--------------------------------------------------------------
%% Repetir el título al inicio del cuerpo (APA 7 §2.11)
\begin{center}
  {\bfseries\docTitulo}
\end{center}

\justifying
\setlength{\parindent}{0.5in}

\section{Introducción}

Antes de la pandemia el Gobierno Boliviano ya habia implementado varias
politicas para impulsar el desarrollo tecnologicos, pero estas medidas solo
estaban pensanda como una herramienta complementaria. Uno de sus proyectos era
la creacion de Telecentros Educativos, donde se proveia de computadoras a
profesores y estudiantes, la informacion se compartia a travez de discos
compactos. Por lo que Bolivia de la misma forma que varios paises no estaba
preparado para el COVID-19.

Este resultado se demuestra a partir de siguiente
\parencite{garcia_zaballos_conectividad_2022} reporte donde muestra estudios que
el servicio de internet en Bolivia es muy caro por que representa
aproximadamente 14\% del salario basico, en comparacion con otros paises que es
el 1 ó 2\%. Asimismo en este estudio se logra visualizar la falta de
insfraestructura con respecto a las tecnologias inalambrias, fibra optica asi
como la cobertura que tenemos no esta a lo que mercado internacional requiere.
En otras palabras los demas paises aumentan su insfraestrcura, y las aplicacion
sitios mobiles y entre otros tambien aumentan la cantidad de infomracion que
manejan. Por lo que nuestra insfraestrcura casi nunca esta a la altura. De la
misma forma este estudio muestra que en muchos lugares de Bolivia aun no llega
la conectividad satelital y entre otros, lo que significa que hay mucha
poblacion especialmente en areas rurales que se encuentran desconectados.

Entonces lo unico que hizo la pandeima fue evidenciar la falta de infraestructura que teniamos

El objetivo principal de este documento es analizar la brecha digital en el
acceso a las TIC en unidades educativas del área rural de Bolivia, identificando
los factores que contribuyen a esta desigualdad y elaborando una propuesta de
solución que permita mejorar la situación actual.

Este documento se basa en la concientizacion de la sociedad sobre las
desigualdades en el acceso a las TIC, y busca generar un cambio positivo en la
educación rural del país. A través de un análisis exhaustivo del problema y la
presentación de una propuesta fundamentada, se espera contribuir al desarrollo
de estrategias que promuevan la equidad en el acceso a las TIC y, por ende, a
una educación de calidad para todos los estudiantes, independientemente de su
ubicación geográfica.

A continuacion el documento se organiza en dos partes. La primera parte se
centra en el planteamiento del problema, mientras que la segunda parte presenta
la propuesta de solución. Finalmente, se incluyen las conclusiones y
recomendaciones del analisis realizado.


\section{Planteamiento del Problema}

\subsection{Análisis del Problema}
\lipsum[1]

\subsection{Fundamentación del Problema}
\lipsum[1]

\section{Propuesta}

\subsection{Esquema de la Propuesta}
\lipsum[1]

\subsection{Descripción de la Propuesta}
\lipsum[1]

\subsection{Fundamentación de la Propuesta}
\lipsum[1]

\section{Conclusión}
\lipsum[1]

\section{Recomendaciones}
\lipsum[8]

\section{Resultados}

La Tabla~\ref{tab:BasicTable} resume los datos obtenidos. \lipsum[15]

\begin{table}[htbp]
  \caption{Tabla básica de ejemplo}
  \label{tab:BasicTable}
  \begin{tabular}{@{}llr@{}}
    \toprule
    \multicolumn{2}{c}{Ítem} \\ \cmidrule(r){1-2}
    Animal    & Descripción  & Precio \\ \midrule
    Gnomo     & por gramo    & 13.65  \\
              & cada uno     &  0.01  \\
    Ñu        & disecado     & 92.50  \\
    Emú       & disecado     & 33.33  \\
    Armadillo & congelado    &  8.99  \\ \bottomrule
  \end{tabular}
\end{table}

\section{Discusión}
\lipsum[17]

%%--------------------------------------------------------------
%% REFERENCIAS
%%--------------------------------------------------------------
\newpage
\printbibliography[title={Referencias Bibliográficas}]

\end{document}
